\documentclass[11pt,a4paper]{article}
\usepackage[utf8]{inputenc}
\usepackage[vietnamese]{babel}
\usepackage{amsmath,amssymb}
\usepackage{geometry}
\geometry{a4paper, total={180mm,260mm}, left=15mm, right=15mm, top=18mm, bottom=18mm}
\usepackage{tikz}
\usetikzlibrary{shapes.geometric, arrows.meta, positioning, calc}
\usepackage{tcolorbox}
\tcbuselibrary{skins,breakable}
\usepackage{xcolor}
\usepackage{multicol}
\usepackage{array}
\usepackage{fontawesome}
\usepackage{fancyhdr}

\pagestyle{fancy}
\fancyhf{}
\renewcommand{\headrulewidth}{0.5pt}
\renewcommand{\footrulewidth}{0.5pt}
\fancyhead[L]{\small\bfseries\color{teal}\faBook\ SỔ TAY KIẾN THỨC TOÁN 3}
\fancyhead[R]{\small\bfseries\color{teal}Lớp Toán Thầy Lĩnh -- loptoanthaylinh.vn}
\fancyfoot[C]{\small\color{gray}-- Sổ tay kiến thức: Bài 30. Mi-li-mét (mm) --}

\definecolor{primary}{HTML}{005F73}
\definecolor{secondary}{HTML}{0A9396}
\definecolor{accent}{HTML}{EE9B00}
\definecolor{danger}{HTML}{AE2012}
\definecolor{boxbg}{HTML}{F4F6F8}

\begin{document}

% --- BANNER TIÊU ĐỀ ---
\begin{tcolorbox}[
  enhanced,
  colback=white,
  colframe=primary,
  arc=3mm,
  boxrule=1.2pt,
  left=4mm, right=4mm, top=2.5mm, bottom=2.5mm,
  drop shadow=black!10!white
]
  \centering
  {\Large\bfseries\color{primary}\faBookmark\ BÀI 30: MI-LI-MÉT (mm)}\\[1.5mm]
  {\small\itshape\color{black!75}Đơn vị đo độ dài siêu nhỏ, bảng chuyển đổi mm, cm, dm, m và thực hành đọc thước kẻ trên 1 trang A4}
\end{tcolorbox}

\vspace{-2mm}


\noindent
\begin{minipage}[t]{0.49\textwidth}\vspace{0pt}
\begin{tcolorbox}[
  enhanced,
  colback=blue!4!white,
  colframe=blue!75!black,
  arc=2mm, boxrule=0.8pt,
  left=2mm, right=2mm, top=1.8mm, bottom=1.8mm,
  title={\footnotesize\bfseries\color{white}\faKey\ 1. Đơn vị Mi-li-mét (mm)}
]
{\footnotesize
\begin{itemize}\setlength{\itemsep}{1.5pt}\setlength{\leftskip}{-2.5mm}
  \item Mi-li-mét là đơn vị đo độ dài, viết tắt là \textbf{mm}.
  \item Dùng để đo những vật có kích thước rất nhỏ hoặc độ dày mỏng (bề dày quyển sách, ngòi bút, hạt gạo...).
  \item \textbf{Quan hệ giữa các đơn vị độ dài:}
  \begin{itemize}\setlength{\itemsep}{1pt}\setlength{\leftskip}{-2.5mm}
    \item $\mathbf{1\text{ cm} = 10\text{ mm}} \iff 10\text{ mm} = 1\text{ cm}$.
    \item $\mathbf{1\text{ dm} = 100\text{ mm}}$.
    \item $\mathbf{1\text{ m} = 1000\text{ mm}}$.
  \end{itemize}
\end{itemize}
}
\end{tcolorbox}
\end{minipage}
\hfill
\begin{minipage}[t]{0.49\textwidth}\vspace{0pt}
\begin{tcolorbox}[
  enhanced,
  colback=teal!4!white,
  colframe=teal!75!black,
  arc=2mm, boxrule=0.8pt,
  left=2mm, right=2mm, top=1.8mm, bottom=1.8mm,
  title={\footnotesize\bfseries\color{white}\faSliders\ 2. Kĩ năng đo bằng thước vạch}
]
{\footnotesize
\begin{itemize}\setlength{\itemsep}{1.5pt}\setlength{\leftskip}{-2.5mm}
  \item Khoảng cách giữa 2 vạch nhỏ liền nhau trên thước kẻ thông thường là $\mathbf{1\text{ mm}}$.
  \item Giữa vạch $0$ và vạch $1\text{ cm}$ có đúng $10$ khoảng nhỏ $\implies 1\text{ cm} = 10\text{ mm}$.
  \item \textbf{Cách đọc độ dài vật:}\\
  Đặt một đầu của vật trùng với vạch số $\mathbf{0}$, đầu kia chỉ vào vạch nào thì đọc số đo theo vạch đó (kèm theo đơn vị).
\end{itemize}
}
\end{tcolorbox}
\end{minipage}

\vspace{-2.5mm}
% --- KHỐI 3: SƠ ĐỒ TƯ DUY TIKZ ---
\noindent
\begin{tcolorbox}[
  enhanced,
  colback=orange!4!white,
  colframe=orange!80!black,
  arc=2mm, boxrule=0.8pt,
  left=2mm, right=2mm, top=1.5mm, bottom=1.5mm,
  title={\footnotesize\bfseries\color{white}\faSitemap\ SƠ ĐỒ TƯ DUY: BẢNG DỊCH CHUYỂN ĐƠN VỊ ĐỘ DÀI}
]
\centering
\begin{tikzpicture}[
  u/.style={draw=blue!80!black, fill=blue!10, rounded corners=2pt, font=\scriptsize\bfseries, minimum width=1.6cm, minimum height=0.65cm, align=center},
  arr/.style={-{Stealth[length=1.8mm]}, thick}
]
  \node[u] (km) at (0, 0) {$\mathbf{km}$};
  \node[u] (m) at (2.6, 0) {$\mathbf{m}$};
  \node[u] (dm) at (5.2, 0) {$\mathbf{dm}$};
  \node[u] (cm) at (7.8, 0) {$\mathbf{cm}$};
  \node[u, fill=orange!20, draw=orange!80!black] (mm) at (10.4, 0) {$\mathbf{mm}$};

  \draw[arr, blue!80!black] (km) -- (m) node[midway, above, font=\tiny\bfseries] {$\times 1000$};
  \draw[arr, teal!80!black] (m) -- (dm) node[midway, above, font=\tiny\bfseries] {$\times 10$};
  \draw[arr, teal!80!black] (dm) -- (cm) node[midway, above, font=\tiny\bfseries] {$\times 10$};
  \draw[arr, orange!80!black] (cm) -- (mm) node[midway, above, font=\tiny\bfseries] {$\times 10$};
\end{tikzpicture}
\end{tcolorbox}

\vspace{-2.5mm}
% --- KHỐI 4: VÍ DỤ & MẸO TRÁNH BẪY ---
\noindent
\begin{minipage}[t]{0.49\textwidth}\vspace{0pt}
\begin{tcolorbox}[
  enhanced,
  colback=green!4!white,
  colframe=green!60!black,
  arc=2mm, boxrule=0.8pt,
  left=2mm, right=2mm, top=1.8mm, bottom=1.8mm,
  title={\footnotesize\bfseries\color{white}\faCheckCircle\ Ví dụ thực hành chuẩn mực}
]
{\footnotesize
\textbf{1. Đổi đơn vị:}\\
$3\text{ cm} = \mathbf{30}\text{ mm}$.\\
$4\text{ cm } 5\text{ mm} = 40\text{ mm} + 5\text{ mm} = \mathbf{45}\text{ mm}$.\\
$70\text{ mm} = 70 : 10 = \mathbf{7}\text{ cm}$.\\
\textbf{2. Thực hiện phép tính có đơn vị:}\\
$25\text{ mm} + 18\text{ mm} = \mathbf{43\text{ mm}}$.\\
$8\text{ mm} \times 4 = \mathbf{32\text{ mm}}$.\\
$36\text{ mm} : 6 = \mathbf{6\text{ mm}}$.
}
\end{tcolorbox}
\end{minipage}
\hfill
\begin{minipage}[t]{0.49\textwidth}\vspace{0pt}
\begin{tcolorbox}[
  enhanced,
  colback=red!4!white,
  colframe=red!70!black,
  arc=2mm, boxrule=0.8pt,
  left=2mm, right=2mm, top=1.8mm, bottom=1.8mm,
  title={\footnotesize\bfseries\color{white}\faExclamationTriangle\ Mẹo tránh bẫy trắc nghiệm}
]
{\footnotesize
\begin{itemize}\setlength{\itemsep}{1.5pt}\setlength{\leftskip}{-2.5mm}
  \item \textbf{Bẫy đặt vật không trùng vạch 0:}\\
  Nếu đầu vật đặt ở vạch $2\text{ cm}$ và đuôi vật ở vạch $5\text{ cm } 4\text{ mm}$, thì độ dài là: $54 - 20 = \mathbf{34\text{ mm}}$, không phải $54\text{ mm}$!
  \item \textbf{Bẫy cộng khác đơn vị:}\\
  Không được cộng trực tiếp $2\text{ cm} + 5\text{ mm} = 7$! Bắt buộc phải đổi về cùng đơn vị: $20\text{ mm} + 5\text{ mm} = \mathbf{25\text{ mm}}$.
\end{itemize}
}
\end{tcolorbox}
\end{minipage}


\end{document}
